\documentclass[a4paper,11pt]{article}
\usepackage{fontspec}
\usepackage{unicode-math}
\usepackage{amsmath,amssymb,amsthm}
\usepackage{longtable,booktabs,array}
\usepackage{hyperref}
\usepackage[margin=25mm]{geometry}
\setmainfont{Latin Modern Roman}

\newtheorem{theorem}{Theorem}[section]
\newtheorem{proposition}[theorem]{Proposition}
\newtheorem{lemma}[theorem]{Lemma}
\newtheorem{corollary}[theorem]{Corollary}
\newtheorem{definition}[theorem]{Definition}
\newtheorem{remark}[theorem]{Remark}
\newtheorem{observation}[theorem]{Observation}

\providecommand{\tightlist}{\setlength{\itemsep}{0pt}\setlength{\parskip}{0pt}}
\newcounter{none}


\begin{document}


\section{D1 v2.0: Particle Data Structure Specification of the Central
Projection Framework --- Frequency Interpretation and Dynamics
Extension}\label{d1-v2.0-particle-data-structure-specification-of-the-central-projection-framework-frequency-interpretation-and-dynamics-extension}

\textbf{Author}: Noriaki Kihara \textbf{Affiliation}: WF System Co.,
Ltd. \textbf{ORCID}: 0009-0004-6753-4020 \textbf{Version}: v2.0
\textbf{Date}: May 2026 \textbf{DOI}:
\href{https://doi.org/10.5281/zenodo.19941250}{10.5281/zenodo.19941250}
\textbf{Predecessor}: v1.0 (April 2026, unpublished)

\begin{center}\rule{0.5\linewidth}{0.5pt}\end{center}

\subsection{§0 Position of This Paper}\label{position-of-this-paper}

This paper is a design specification, not a mathematical proof. It
explicitly defines, as a typed-integer-arithmetic specification, the
particle data structures that have been implicitly assumed in the
preceding paper series {[}W1{]}--{[}W11{]}, {[}M1{]}--{[}M3{]},
supplementary lectures {[}Sup1{]}--{[}Sup4{]}, and the thought
experiment {[}TE\_Raxis{]}.

\subsubsection{Major Revisions from
v1.0}\label{major-revisions-from-v1.0}

\begin{enumerate}
\def\labelenumi{\arabic{enumi}.}
\tightlist
\item
  \textbf{Reinterpret σ.k\_xyzt as ``vibrational frequencies of each
  axis''} rather than ``wave numbers / acceleration'' (§1.4)
\item
  \textbf{Add δV (phase displacement) and δδt (processing-period change)
  to the Particle structure} (§5)
\item
  \textbf{Introduce the Mesh structure to separate longitudinal modes
  (Higgs, gravitational waves) from Particles} (§7)
\item
  \textbf{Eliminate special status of t-axis, fully preserving 6-axis
  symmetry} (§3, §6)
\end{enumerate}

The data type system (int256-based) and integer-arithmetic principles
(no division, no square root, no trigonometric functions) remain
unchanged from v1.0.

\subsubsection{Claims of This Paper}\label{claims-of-this-paper}

\begin{enumerate}
\def\labelenumi{\arabic{enumi}.}
\tightlist
\item
  The boundedness of the universe and the discreteness at the Planck
  scale together permit all phase quantities to be represented as
  256-bit signed integers (§1).
\item
  The \([L^{-1}]\) dimension of the xyzt axes is naturally interpreted
  as ``vibrational frequency along each axis in phase space,'' fully
  preserving 6-axis symmetry (§3).
\item
  The complete state of a particle can be described by 5 fields:
  \texttt{SignVector} (identity label) + \texttt{OneDWave{[}κ{]}} (phase
  state) + \texttt{δV} (phase displacement) + \texttt{δt,\ δδt}
  (processing period and acceleration) (§5).
\item
  Longitudinal modes such as the Higgs should be represented not as
  Particles but as vibrations of the \texttt{Mesh}, allowing unified
  treatment of mass generation, gravitational waves, and cosmological
  vibrations (§7).
\item
  All operations close under integer addition and multiplication,
  requiring no division, square root, or trigonometric functions (§6).
\end{enumerate}

\subsubsection{Out of Scope}\label{out-of-scope}

\begin{itemize}
\tightlist
\item
  Dynamical formulation of operation rules (physical meaning of phase
  evolution, concrete forms of interactions) --- deferred to a separate
  paper (D2).
\item
  Numerical derivation of specific physical quantities.
\item
  Hardware implementation details.
\item
  Mathematical proofs of theorems.
\end{itemize}

\subsubsection{Referenced Papers}\label{referenced-papers}

\begin{itemize}
\tightlist
\item
  {[}W3{]} Structural Requirements for a Theory of Everything (DOI:
  10.5281/zenodo.19601592)
\item
  {[}W8{]} Combinatorial Structure of the 6-Dimensional Hyperrectangle
  (DOI: 10.5281/zenodo.19762134)
\item
  {[}W9{]} sine-Gordon Equation --- Foundation of Topological Solitons
  (DOI: 10.5281/zenodo.19650966)
\item
  {[}W10{]} Four Modes of Shape-Invariant Waves (DOI:
  10.5281/zenodo.19709798)
\item
  {[}W11{]} Interactions of Shape-Invariant Waves (DOI:
  10.5281/zenodo.19763463)
\item
  {[}M1{]} Three-Layer Model of Central Projection: R--R₁--R₀ (DOI:
  10.5281/zenodo.19691713)
\item
  {[}TE\_Raxis{]} Reexamination of the 6-Dimensional Sign Vector xyztRQ
  (DOI: 10.5281/zenodo.19904714)
\end{itemize}

\begin{center}\rule{0.5\linewidth}{0.5pt}\end{center}

\subsection{§1 int256 Boundedness and Dimensional
Interpretation}\label{int256-boundedness-and-dimensional-interpretation}

\subsubsection{§1.1 Basis for Boundedness Requirement (unchanged from
v1.0)}\label{basis-for-boundedness-requirement-unchanged-from-v1.0}

{[}W3{]} requires boundedness as a structural condition for any theory
of everything. In the three-layer structure \(R\)--\(R_1\)--\(R_0\) of
{[}M1{]}, the maximum curvature radius \(R_0\) is finite. Combined with
discreteness at the Planck scale, this yields an upper bound on the bits
needed to represent all phase quantities.

\subsubsection{§1.2 Required Number of
Bits}\label{required-number-of-bits}

\textbf{Spatial axis}: From the observable universe diameter
\(D \approx 1.892 \times 10^{27}\) m and Planck length
\(l_P = 1.616 \times 10^{-35}\) m,

\[n_{\text{space}} = D / l_P \approx 1.171 \times 10^{62}, \quad b_{\text{space}} = 207 \text{ bits}\]

\textbf{Time axis}: From cosmic time \(T \approx 8.710 \times 10^{17}\)
s and Planck time \(t_P = 5.391 \times 10^{-44}\) s,

\[n_{\text{time}} = T / t_P \approx 1.616 \times 10^{61}, \quad b_{\text{time}} = 204 \text{ bits}\]

\textbf{Conclusion}: 256-bit signed integers leave 14-digit margin in
space and 15-digit margin in time.

\subsubsection{§1.3 Natural Units and
Dimensions}\label{natural-units-and-dimensions}

Under \(c = 1\), \(G = 1\) (geometrized units; \(\hbar\) is not
adopted),

\[[L] = [T] = [M], \quad [L^{-1}] = [M^{-1}]\]

\subsubsection{§1.4 Frequency Interpretation (Core Revision in
v2.0)}\label{frequency-interpretation-core-revision-in-v2.0}

In v1.0, the \([L^{-1}]\) dimension of the xyzt axes was interpreted
ambiguously as ``wave number'' or ``acceleration.'' In v2.0, this is
unified as \textbf{``vibrational frequency along each axis''} as the
natural interpretation in phase space.

{\def\LTcaptype{none} % do not increment counter
\begin{longtable}[]{@{}
  >{\raggedright\arraybackslash}p{(\linewidth - 6\tabcolsep) * \real{0.2500}}
  >{\raggedright\arraybackslash}p{(\linewidth - 6\tabcolsep) * \real{0.2500}}
  >{\raggedright\arraybackslash}p{(\linewidth - 6\tabcolsep) * \real{0.2500}}
  >{\raggedright\arraybackslash}p{(\linewidth - 6\tabcolsep) * \real{0.2500}}@{}}
\toprule\noalign{}
\begin{minipage}[b]{\linewidth}\raggedright
Quantity
\end{minipage} & \begin{minipage}[b]{\linewidth}\raggedright
Dimension
\end{minipage} & \begin{minipage}[b]{\linewidth}\raggedright
v1.0 Interpretation
\end{minipage} & \begin{minipage}[b]{\linewidth}\raggedright
\textbf{v2.0 Interpretation}
\end{minipage} \\
\midrule\noalign{}
\endhead
\bottomrule\noalign{}
\endlastfoot
σ.k\_x, k\_y, k\_z & \([L^{-1}]\) & spatial wave number / acceleration &
\textbf{vibrational frequency of each spatial axis} \\
σ.k\_t & \([L^{-1}]\) & temporal frequency & \textbf{vibrational
frequency of t-axis} (equal to others) \\
σ.R & \([L]=[M]\) & curvature radius & \textbf{scale quantity /
corresponding period} (dual) \\
σ.Q & \([1]\) & internal symmetry & internal symmetry (unchanged) \\
δt & \([L]=[M]\) & proper time step & \textbf{processing step period}
(not ``time'') \\
δδt & \([L]=[M]\) & (not in v1.0) & \textbf{change in processing period}
(acceleration effect) \\
\end{longtable}
}

\subsubsection{§1.5 Advantages of the Frequency
Interpretation}\label{advantages-of-the-frequency-interpretation}

\textbf{(1) Complete preservation of 6-axis symmetry}: All four
position-type axes (xyzt) carry independent vibrational modes; the
t-axis is on equal footing with the others.

\textbf{(2) Elimination of continuum-mechanical concepts}:
``Acceleration'' is a continuum concept (second time derivative of
position) that does not fit naturally into discrete integer arithmetic.
``Vibrational frequency'' is the natural discrete quantity in phase
space.

\textbf{(3) Clear separation of acceleration effects}: Acceleration-like
effects (gravitational fields, inertial frame influences) are
represented by the change δδt of the processing period, structurally
separated from the particle's intrinsic frequency σ.k.

\textbf{(4) Frequency × period = phase}: \(\sigma.k \cdot \delta t\)
becomes a dimensionless phase increment (\([L^{-1}] \cdot [L] = [1]\)),
making the phase evolution rule type-consistent.

\textbf{(5) Automatic inclusion of special relativity}: The dispersion
relation
\(\sigma.k_t^2 = \sigma.k_x^2 + \sigma.k_y^2 + \sigma.k_z^2 + m^2\) can
be naturally written as integer relations among frequencies.

\begin{center}\rule{0.5\linewidth}{0.5pt}\end{center}

\subsection{§2 Definition of Meta Types}\label{definition-of-meta-types}

\subsubsection{§2.1 Scalar Types (unchanged from
v1.0)}\label{scalar-types-unchanged-from-v1.0}

\begin{verbatim}
Type         Dimension  Range                     Use
L256         [L]=[M]    |v| ≤ 2^255 - 1          length / time / mass / period
Linv256      [L⁻¹]      |v| ≤ 2^255 - 1          frequency / inverse scale
Scalar256    [1]        |v| ≤ 2^255 - 1          dimensionless
FullPhase    [1]        0 ≤ v ≤ P                0–720°
CyclePhase   [1]        0 ≤ v ≤ P/2              0–360° (position phase)
DeltaPhase   [1]        -P/2 ≤ v ≤ P/2           ±360° (phase displacement)
SpinPhase    [1]        v ∈ {0, P/4, P/2, P}     {0, 180, 360, 720}°
\end{verbatim}

Phase constant \(P = \lfloor (2^{255}-1)/720 \rfloor \times 720\).

\subsubsection{§2.2 Axis Component Types (renamed and reinterpreted in
v2.0)}\label{axis-component-types-renamed-and-reinterpreted-in-v2.0}

{\def\LTcaptype{none} % do not increment counter
\begin{longtable}[]{@{}
  >{\raggedright\arraybackslash}p{(\linewidth - 6\tabcolsep) * \real{0.2500}}
  >{\raggedright\arraybackslash}p{(\linewidth - 6\tabcolsep) * \real{0.2500}}
  >{\raggedright\arraybackslash}p{(\linewidth - 6\tabcolsep) * \real{0.2500}}
  >{\raggedright\arraybackslash}p{(\linewidth - 6\tabcolsep) * \real{0.2500}}@{}}
\toprule\noalign{}
\begin{minipage}[b]{\linewidth}\raggedright
Type
\end{minipage} & \begin{minipage}[b]{\linewidth}\raggedright
Base
\end{minipage} & \begin{minipage}[b]{\linewidth}\raggedright
Dimension
\end{minipage} & \begin{minipage}[b]{\linewidth}\raggedright
Physical Meaning (v2.0)
\end{minipage} \\
\midrule\noalign{}
\endhead
\bottomrule\noalign{}
\endlastfoot
\texttt{SpatialFreq} \textless: Linv256 & int256 & \([L^{-1}]\) &
vibrational frequency of x, y, z axes \\
\texttt{TemporalFreq} \textless: Linv256 & int256 & \([L^{-1}]\) &
vibrational frequency of t-axis (equal to others) \\
\texttt{CurvatureR} \textless: L256 & int256 & \([L]\) & scale / period
(signed square number) \\
\texttt{ColorQ} \textless: Scalar256 & int256 & \([1]\) & internal
symmetry label (finite group) \\
\texttt{DeltaT} \textless: L256 & int256 & \([L]\) & processing step
period \\
\end{longtable}
}

\textbf{Renaming from v1.0}: \texttt{SpatialK} → \texttt{SpatialFreq},
\texttt{TemporalK} → \texttt{TemporalFreq}. The type system itself is
unchanged; only the interpretation is clarified.

\subsubsection{§2.3 Operation Rules (unchanged from
v1.0)}\label{operation-rules-unchanged-from-v1.0}

\textbf{Addition}: Only between same types (CyclePhase + DeltaPhase =
CyclePhase, etc.).

\textbf{Multiplication}: - \texttt{L256\ ×\ L256\ →\ L2\_256} \([L^2]\)
- \texttt{Linv256\ ×\ Linv256\ →\ Linv2\_256} \([L^{-2}]\) -
\texttt{Linv256\ ×\ L256\ →\ Scalar256} \([1]\) (dimensionless)

\textbf{Division}: Not used.

\begin{center}\rule{0.5\linewidth}{0.5pt}\end{center}

\subsection{§3 Definition of SignVector (Frequency Interpretation
Version)}\label{definition-of-signvector-frequency-interpretation-version}

\subsubsection{§3.1 Structure}\label{structure}

\begin{verbatim}
SignVector = (
    x : SpatialFreq,    -- frequency of x-axis vibration [L⁻¹]
    y : SpatialFreq,    -- frequency of y-axis vibration
    z : SpatialFreq,    -- frequency of z-axis vibration
    t : TemporalFreq,   -- frequency of t-axis vibration [L⁻¹] (equal to others)
    R : CurvatureR,     -- scale quantity / period [L]
    Q : ColorQ,         -- internal symmetry label [1]
)
\end{verbatim}

\subsubsection{§3.2 Physical Meaning}\label{physical-meaning}

Each σ component represents either a ``vibrational frequency along the
corresponding axis'' or a ``scale quantity'':

\begin{itemize}
\tightlist
\item
  σ.k\_axis = 0: the axis does not vibrate (no structure in that
  direction)
\item
  σ.k\_axis ≠ 0: the axis vibrates at frequency
  \textbar k\_axis\textbar{}
\item
  sign(k\_axis): direction of phase rotation (+ for clockwise, − for
  counterclockwise)
\item
  \textbar k\_axis\textbar: magnitude of vibrational frequency (discrete
  integer value)
\end{itemize}

\subsubsection{§3.3 Invariants}\label{invariants}

The values of SignVector change in interaction events, but specific
combinations define stable modes:

\begin{itemize}
\tightlist
\item
  \(\kappa_s = |\{i \in \{x,y,z\} : k_i \neq 0\}|\) --- number of
  non-zero spatial axes
\item
  \(\kappa_t = (1 \text{ if } k_t \neq 0 \text{ else } 0)\) --- non-zero
  indicator for t-axis
\item
  \(\kappa = \kappa_s + \kappa_t \in \{0, 1, 2\}\) --- {[}W10{]}
  Proposition 2.1: \(\kappa \geq 3\) is unstable
\end{itemize}

\subsubsection{§3.4 Dispersion Relation}\label{dispersion-relation}

The mass-shell condition is directly expressed as an integer relation
among frequencies:

\[\sigma.k_t^2 = \sigma.k_x^2 + \sigma.k_y^2 + \sigma.k_z^2 + m^2\]

\begin{itemize}
\tightlist
\item
  \(m = 0\) (photon, gluon): \(k_t^2 = k_x^2 + k_y^2 + k_z^2\)
\item
  \(m \neq 0\) (massive particle): the square root of the difference
  corresponds to mass
\end{itemize}

This automatically incorporates special relativity into the framework.

\subsubsection{§3.5 Correspondence with {[}W8{]} Table A (continued from
v1.0)}\label{correspondence-with-w8-table-a-continued-from-v1.0}

The sign patterns \(\{+, -, 0\}^6\) in {[}W8{]} Table A correspond to
the signs of SignVector. Integer values (axis scale values) correspond
to \(v_i\) in {[}Mass Structure{]} §1.

Representative particles:

{\def\LTcaptype{none} % do not increment counter
\begin{longtable}[]{@{}
  >{\raggedright\arraybackslash}p{(\linewidth - 6\tabcolsep) * \real{0.2500}}
  >{\raggedright\arraybackslash}p{(\linewidth - 6\tabcolsep) * \real{0.2500}}
  >{\raggedright\arraybackslash}p{(\linewidth - 6\tabcolsep) * \real{0.2500}}
  >{\raggedright\arraybackslash}p{(\linewidth - 6\tabcolsep) * \real{0.2500}}@{}}
\toprule\noalign{}
\begin{minipage}[b]{\linewidth}\raggedright
Particle
\end{minipage} & \begin{minipage}[b]{\linewidth}\raggedright
Set \((x,y,z,t,R,Q)\)
\end{minipage} & \begin{minipage}[b]{\linewidth}\raggedright
\(\kappa\)
\end{minipage} & \begin{minipage}[b]{\linewidth}\raggedright
Classification
\end{minipage} \\
\midrule\noalign{}
\endhead
\bottomrule\noalign{}
\endlastfoot
Photon \(\gamma\) & \((0,0,0,0,+,0)\) & 0 & gauge boson (free, spherical
wave) \\
Gluon \(g\) & \((0,0,0,0,0,Q_i \to Q_j)\) & 0 & gauge boson \\
Higgs \(H\) & \((0,0,+,0,0,0)\) & 1 & scalar (longitudinal
interpretation in §7) \\
\(W^+\) & \((+,0,0,0,+,0)\) & 1 & gauge boson \\
\(Z^0\) & \((0,+,0,0,+,0)\) & 1 & gauge boson \\
Graviton & \((0,0,0,\pm,\pm,0)\) & 1 & tensor boson \\
Electron \(e^-\) & \((+,0,0,+,0,0)\) & 2 & fermion \\
\end{longtable}
}

\begin{center}\rule{0.5\linewidth}{0.5pt}\end{center}

\subsection{§4 Definition of OneDWave}\label{definition-of-onedwave}

\subsubsection{§4.1 Structure (unchanged from
v1.0)}\label{structure-unchanged-from-v1.0}

\begin{verbatim}
OneDWave = (
    θ₀ : CyclePhase,   -- position phase [1]   (S1 §2.1, W9 §8.2)
    A  : Scalar256,     -- amplitude [1]        (S1 §2.1, W9 §8.2)
    ℓ  : DeltaPhase,    -- spread phase [1]     (W9 §8.3)
)
\end{verbatim}

\subsubsection{§4.2 Physical Meaning}\label{physical-meaning-1}

Each particle has one OneDWave per non-zero position-type axis (at most
\(\kappa\)). Each OneDWave holds the current phase state of vibration in
that axis direction.

\begin{itemize}
\tightlist
\item
  \(\theta_0\): current phase of vibration (free parameter, evolves over
  time)
\item
  \(A\): amplitude (free parameter)
\item
  \(\ell\): spatial extent of localization (dependent on SignVector,
  {[}W9{]} §8.3)
\end{itemize}

\begin{center}\rule{0.5\linewidth}{0.5pt}\end{center}

\subsection{§5 Definition of Particle (δV and δδt added in
v2.0)}\label{definition-of-particle-ux3b4v-and-ux3b4ux3b4t-added-in-v2.0}

\subsubsection{§5.1 Structure}\label{structure-1}

\begin{verbatim}
Particle = (
    σ   : SignVector,    -- identity (vibrational frequencies of each axis)
    V   : OneDWave[κ],   -- current phase state of each non-zero position-type axis
    δV  : OneDWave[κ],   -- phase displacement per processing step (added in v2.0)
    δt  : DeltaT,         -- current processing step period
    δδt : DeltaT,         -- change in processing period (added in v2.0, ≈ acceleration)
)
\end{verbatim}

\subsubsection{§5.2 Physical Meaning of Each
Field}\label{physical-meaning-of-each-field}

\paragraph{\texorpdfstring{\texttt{σ\ :\ SignVector}}{σ : SignVector}}\label{ux3c3-signvector}

The particle's identity. Specifies frequencies of vibration in each
axis. Values change in interaction events (formulated as annihilation +
creation, but as a data structure σ is updated).

\paragraph{\texorpdfstring{\texttt{V\ :\ OneDWave{[}κ{]}}}{V : OneDWave{[}κ{]}}}\label{v-onedwaveux3ba}

Current phase state of each non-zero position-type axis. Array length
\(\kappa\) equals the number of non-zero position-type axes in
SignVector (\(\leq 2\)). Each V{[}j{]} holds the (position phase,
amplitude, spread) of vibration in the corresponding axis.

Index correspondence:

\begin{verbatim}
σ.k_x ≠ 0  ⇒  V contains wave_x
σ.k_y ≠ 0  ⇒  V contains wave_y
σ.k_z ≠ 0  ⇒  V contains wave_z
σ.k_t ≠ 0  ⇒  V contains wave_t
\end{verbatim}

\paragraph{\texorpdfstring{\texttt{δV\ :\ OneDWave{[}κ{]}} (added in
v2.0)}{δV : OneDWave{[}κ{]} (added in v2.0)}}\label{ux3b4v-onedwaveux3ba-added-in-v2.0}

A difference structure isomorphic to V. Holds how V changes in one
processing step:

\begin{verbatim}
δV[j] = (
    Δθ₀ : DeltaPhase,    -- phase displacement
    ΔA  : Scalar256,      -- amplitude displacement
    Δℓ  : DeltaPhase,     -- spread displacement
)
\end{verbatim}

For free particles, δV{[}j{]}.Δθ₀ can be computed from σ:

\[\delta V[j].\Delta\theta_0 = \sigma.k_{a(j)} \cdot \delta t\]

However, in interaction events δV may change independently.

\paragraph{\texorpdfstring{\texttt{δt\ :\ DeltaT}}{δt : DeltaT}}\label{ux3b4t-deltat}

\textbf{Processing step period} (NOT time). A universal computation
counter.

\begin{itemize}
\tightlist
\item
  \(\delta t > 0\): forward operation (causal forward direction)
\item
  \(\delta t < 0\): reverse operation (guarantees reversibility)
\item
  \(\delta t = 0\): at rest
\end{itemize}

What is observed as ``time'' is the accumulation of V{[}t-axis{]}.θ₀,
not δt itself (see §6.4).

\paragraph{\texorpdfstring{\texttt{δδt\ :\ DeltaT} (added in
v2.0)}{δδt : DeltaT (added in v2.0)}}\label{ux3b4ux3b4t-deltat-added-in-v2.0}

\textbf{Change in processing period} (acceleration equivalent).
Determined by background R(r) or interactions:

\begin{verbatim}
1-step update: δt += δδt
\end{verbatim}

Physical interpretation:

\begin{itemize}
\tightlist
\item
  \(\delta\delta t = 0\): free particle (constant processing period)
\item
  \(\delta\delta t > 0\): processing period elongates (clock slows down
  in gravity / acceleration)
\item
  \(\delta\delta t < 0\): processing period shortens (rising
  gravitational potential, etc.)
\end{itemize}

\subsubsection{§5.3 Hierarchy Closes at
δδt}\label{hierarchy-closes-at-ux3b4ux3b4t}

Because the basic laws of physics close at second order in time
(Newtonian mechanics, GR, quantum mechanics), no higher-order structure
(δδδt, etc.) is needed. \textbf{The 5 fields (σ, V, δV, δt, δδt) form a
complete representation for dynamics.}

\subsubsection{§5.4 Memory Footprint}\label{memory-footprint}

\begin{verbatim}
σ      : 6 × 256 = 1,536 bit
V      : κ × 3 × 256 ≤ 1,536 bit (κ ≤ 2)
δV     : κ × 3 × 256 ≤ 1,536 bit
δt     : 256 bit
δδt    : 256 bit
Total  ≤ 5,120 bit = 640 byte / particle
\end{verbatim}

This is increased from v1.0's 224--416 byte, but still extremely small
for a complete description of any particle in the universe.

\begin{center}\rule{0.5\linewidth}{0.5pt}\end{center}

\subsection{§6 Operation Rules}\label{operation-rules}

\subsubsection{§6.1 Basic: 1-Step Update}\label{basic-1-step-update}

For each non-zero position-type axis \(j\):

\[\delta V[j].\Delta\theta_0 \leftarrow \sigma.k_{a(j)} \cdot \delta t\]

\[V[j].\theta_0 \mathrel{+}= \delta V[j].\Delta\theta_0\]

Then update the processing period itself:

\[\delta t \mathrel{+}= \delta\delta t\]

\subsubsection{§6.2 Type Consistency
Check}\label{type-consistency-check}

\begin{verbatim}
σ.k_axis · δt : SpatialFreq × DeltaT 
              = Linv256 × L256 
              = Scalar256 [1]                ✓

V[j].θ₀ + δV[j].Δθ₀ : CyclePhase + DeltaPhase 
                    = CyclePhase [1] (mod P/2)  ✓

δt + δδt : DeltaT + DeltaT 
         = DeltaT [L]                       ✓
\end{verbatim}

\subsubsection{§6.3 Polarization and Multi-Axis Phase
Difference}\label{polarization-and-multi-axis-phase-difference}

For particles with multiple non-zero axes (e.g., a polarized directional
photon):

\begin{verbatim}
σ_polarized_light = (+k_x, +k_y, 0, 0, +R_γ, 0)
V = [wave_x, wave_y]
\end{verbatim}

Polarization phase difference:

\[\Delta\varphi = V[0].\theta_0 - V[1].\theta_0 \in \text{CyclePhase}\]

{\def\LTcaptype{none} % do not increment counter
\begin{longtable}[]{@{}ll@{}}
\toprule\noalign{}
Polarization State & \(\Delta\varphi\) \\
\midrule\noalign{}
\endhead
\bottomrule\noalign{}
\endlastfoot
Linear (in-phase) & \(0\) \\
Linear (anti-phase) & \(P/4 = 180°\) \\
Circular (left) & \(+P/8 = +90°\) \\
Circular (right) & \(-P/8 = -90°\) \\
Elliptical & any integer value \\
\end{longtable}
}

All exactly representable as int256 integers (integer multiples of
\(15°\) correspond to integer multiples of \(P/48\) without error).

\subsubsection{§6.4 Treatment of Time: δt is a Processing Step, Time is
V{[}t{]}.θ₀}\label{treatment-of-time-ux3b4t-is-a-processing-step-time-is-vt.ux3b8ux2080}

\textbf{Important principle}: δt is not time but a universal processing
step. What is observed as time is the accumulation of V{[}t-axis{]}.θ₀
for each particle.

Electron example (σ.k\_t ≠ 0):

\begin{verbatim}
σ_electron = (+k_x, 0, 0, +k_t, 0, 0)
V = [wave_x, wave_t]
1-step:
  δV[0].Δθ₀ = σ.k_x · δt    -- x-axis phase displacement (momentum equivalent)
  δV[1].Δθ₀ = σ.k_t · δt    -- t-axis phase displacement (electron's proper time progress)
  V[0].θ₀ += δV[0].Δθ₀
  V[1].θ₀ += δV[1].Δθ₀       -- "time experienced by electron"
  δt += δδt                    -- update processing period (δδt ≠ 0 in gravity)
\end{verbatim}

Photon example (σ.k\_t = 0):

\begin{verbatim}
σ_photon = (+k_x, +k_y, 0, 0, +R_γ, 0)
V = [wave_x, wave_y]      -- no t-axis component
δV = [δwave_x, δwave_y]
-- because σ.k_t = 0, no t-axis OneDWave exists in V
-- consequence: zero proper time of photon is structurally guaranteed
\end{verbatim}

\subsubsection{§6.5 Acceleration Effect (Background
Coupling)}\label{acceleration-effect-background-coupling}

δδt is determined by the background R(r) or by interactions. The
concrete decision rule is defined in a separate paper (D2).

Qualitatively:

\begin{itemize}
\tightlist
\item
  In a gravitational field:
  \(\delta\delta t = f(R_{\text{background}}, \sigma)\) --- background
  curvature modulates processing period
\item
  Result: σ.k · δt changes → frequency shift / time delay / polarization
  phase rotation
\end{itemize}

All of these are expressed as phase evolution by integer arithmetic.

\subsubsection{§6.6 Operations Not Required (same as
v1.0)}\label{operations-not-required-same-as-v1.0}

The following operations are not required in this specification:

\begin{enumerate}
\def\labelenumi{\arabic{enumi}.}
\tightlist
\item
  \textbf{Division}: \(R/R_1\) is converted to \(R \cdot R_1^{-1}\) as
  multiplication
\item
  \textbf{Square root}: a continuum concept; in discrete systems,
  integer values directly determine physical quantities
\item
  \textbf{Trigonometric functions}: tools of continuum sine-Gordon;
  replaceable by discrete modulation tables
\item
  \textbf{Floating-point arithmetic}: all fields are int256, no rounding
  error structurally arises
\end{enumerate}

Where cos/sin would be needed (e.g., longitudinal mode modulation),
discrete modulation tables (lookup tables) of finite resolution are
used.

\begin{center}\rule{0.5\linewidth}{0.5pt}\end{center}

\subsection{§7 Introduction of the Mesh Structure (added in
v2.0)}\label{introduction-of-the-mesh-structure-added-in-v2.0}

\subsubsection{§7.1 Separation of Longitudinal
Modes}\label{separation-of-longitudinal-modes}

In the interpretation of {[}W8{]} Table A, the Higgs
(\(\sigma_H = (0, 0, +k_z, 0, 0, 0)\)) differs from other spin-1 bosons
in having \(\sigma.R = 0\) and \(\text{spin} = 0\). Structurally it is
natural to treat the Higgs not as a ``transverse particle vibrating
position phase'' but as a \textbf{``longitudinal vibration of the mesh
itself.''}

Similarly, the graviton (gravitational waves) can be interpreted as
longitudinal-like vibration of the spacetime metric --- a vibrational
mode of the Mesh.

Treating these longitudinal modes as Particles strains the data
structure (OneDWave is essentially a type for phase rotation and cannot
directly represent axial expansion-contraction). This paper introduces a
\textbf{Mesh structure to manage longitudinal modes separately from
Particles}.

\subsubsection{§7.2 Structure of Mesh}\label{structure-of-mesh}

\begin{verbatim}
Mesh = (
    -- background scale
    R_0           : L256,             -- central length of central projection (background)
    
    -- longitudinal modes
    H_xyz_inphase : OneDWave,         -- xyz in-phase expansion-contraction = Higgs
    H_xy_anti     : OneDWave,         -- xy anti-phase = gravitational wave + polarization
    H_xy_quarter  : OneDWave,         -- xy 90° phase = gravitational wave × polarization
    H_yz_anti     : OneDWave,         -- yz anti-phase = other gravitational wave polarizations
    H_xz_anti     : OneDWave,         -- xz anti-phase = other gravitational wave polarizations
    
    -- displacements of longitudinal modes
    δH_xyz_inphase : OneDWave,
    δH_xy_anti     : OneDWave,
    δH_xy_quarter  : OneDWave,
    δH_yz_anti     : OneDWave,
    δH_xz_anti     : OneDWave,
    
    -- global processing step
    δt_mesh       : DeltaT,
    δδt_mesh      : DeltaT,
)
\end{verbatim}

\subsubsection{§7.3 Physical Identification of Longitudinal
Modes}\label{physical-identification-of-longitudinal-modes}

{\def\LTcaptype{none} % do not increment counter
\begin{longtable}[]{@{}
  >{\raggedright\arraybackslash}p{(\linewidth - 6\tabcolsep) * \real{0.2500}}
  >{\raggedright\arraybackslash}p{(\linewidth - 6\tabcolsep) * \real{0.2500}}
  >{\raggedright\arraybackslash}p{(\linewidth - 6\tabcolsep) * \real{0.2500}}
  >{\raggedright\arraybackslash}p{(\linewidth - 6\tabcolsep) * \real{0.2500}}@{}}
\toprule\noalign{}
\begin{minipage}[b]{\linewidth}\raggedright
Mode
\end{minipage} & \begin{minipage}[b]{\linewidth}\raggedright
Vibration Pattern
\end{minipage} & \begin{minipage}[b]{\linewidth}\raggedright
Physical Identification
\end{minipage} & \begin{minipage}[b]{\linewidth}\raggedright
Spin
\end{minipage} \\
\midrule\noalign{}
\endhead
\bottomrule\noalign{}
\endlastfoot
H\_xyz\_inphase & xyz axes in-phase expansion-contraction & Higgs
(scalar) & 0 \\
H\_xy\_anti & x stretches while y compresses (anti-phase) &
gravitational wave + polarization & 2 \\
H\_xy\_quarter & xy 90° phase rotation & gravitational wave ×
polarization & 2 \\
H\_yz\_anti & yz anti-phase & gravitational wave (other polarization) &
2 \\
H\_xz\_anti & xz anti-phase & gravitational wave (other polarization) &
2 \\
\end{longtable}
}

\textbf{Higgs = spherically symmetric in-phase longitudinal vibration of
xyz axes}: the choice of z-axis in W8 is a representative notation; the
essence is a 3-axis symmetric longitudinal vibration.

\textbf{Two polarizations of gravitational waves (+, ×)}: correspond to
H\_xy\_anti and H\_xy\_quarter (and other axis selections give other
polarizations).

\subsubsection{§7.4 Coupling between Mesh and
Particle}\label{coupling-between-mesh-and-particle}

The vibrational frequency σ.k of a particle is effectively modulated by
the state of the Mesh. Mass generation and frequency shifts are
expressed as this modulation.

Effective frequency:

\[\sigma.k_{\text{effective}}(particle, mesh) = \sigma.k(particle) \cdot \text{mod}(mesh)\]

\[\text{mod}(mesh) = 1 + \alpha_H \cdot \text{disc-cos}(H_{\text{xyz\_inphase}}.\theta_0) + \cdots\]

Here \texttt{disc-cos} is a finite-resolution discrete modulation table
(e.g., 12 stages at 30° intervals) substituting for the continuous
trigonometric function.

\subsubsection{§7.5 Dynamics}\label{dynamics}

Mesh update:

\begin{verbatim}
For each longitudinal mode H_*:
    δH_*.Δθ₀ ← (intrinsic frequency) · δt_mesh
    H_*.θ₀  += δH_*.Δθ₀
δt_mesh += δδt_mesh
\end{verbatim}

Particle update (subordinate to Mesh):

\begin{verbatim}
For each particle p:
    For each non-zero axis j:
        δV[j].Δθ₀ ← σ.k_effective(p, mesh) · δt
        V[j].θ₀  += δV[j].Δθ₀
    δt += δδt
\end{verbatim}

\subsubsection{§7.6 Memory Footprint}\label{memory-footprint-1}

\begin{verbatim}
R_0    : 256 bit
H_*    : 5 × 3 × 256 = 3,840 bit
δH_*   : 5 × 3 × 256 = 3,840 bit
δt_mesh, δδt_mesh : 2 × 256 = 512 bit
Total  = 8,448 bit = 1,056 byte / entire universe
\end{verbatim}

About 1 KB for the entire universe. For \(N\) particles, total memory is
\(640N + 1056\) byte.

\begin{center}\rule{0.5\linewidth}{0.5pt}\end{center}

\subsection{§8 Verification of W3 Structural Requirements (updated in
v2.0)}\label{verification-of-w3-structural-requirements-updated-in-v2.0}

{\def\LTcaptype{none} % do not increment counter
\begin{longtable}[]{@{}
  >{\raggedright\arraybackslash}p{(\linewidth - 4\tabcolsep) * \real{0.3333}}
  >{\raggedright\arraybackslash}p{(\linewidth - 4\tabcolsep) * \real{0.3333}}
  >{\raggedright\arraybackslash}p{(\linewidth - 4\tabcolsep) * \real{0.3333}}@{}}
\toprule\noalign{}
\begin{minipage}[b]{\linewidth}\raggedright
W3 Requirement
\end{minipage} & \begin{minipage}[b]{\linewidth}\raggedright
Satisfied
\end{minipage} & \begin{minipage}[b]{\linewidth}\raggedright
Basis
\end{minipage} \\
\midrule\noalign{}
\endhead
\bottomrule\noalign{}
\endlastfoot
Boundedness & ✓ & int256 represents all phase quantities (§1) \\
Discreteness & ✓ & all fields are integer types (§2--§5) \\
Information conservation & ✓ & δt \textless{} 0 enables reverse
operation (§5.2) \\
Avoidance of zero division & ✓ & division not used (§6.6) \\
Closure under integer arithmetic & ✓ & only addition/multiplication; cos
via discrete modulation table (§7.4) \\
\textbf{6-axis symmetry} & \textbf{✓ (v2.0)} & frequency interpretation
makes all axes equivalent (§1.4, §3) \\
\textbf{Transverse-longitudinal separation} & \textbf{✓ (v2.0)} &
Particle vs Mesh (§7) \\
\textbf{Natural representation of acceleration} & \textbf{✓ (v2.0)} &
δδt is explicit (§5.2) \\
Non-privileging of time & \textbf{✓ (v2.0)} & t equal to other axes;
time is V{[}t{]}.θ₀ (§6.4) \\
Finite representability & ✓ & ≤ 640 byte per particle, +1 KB for whole
universe (§5.4, §7.6) \\
\end{longtable}
}

\begin{center}\rule{0.5\linewidth}{0.5pt}\end{center}

\subsection{§9 Incompatibilities with
v1.0}\label{incompatibilities-with-v1.0}

\subsubsection{§9.1 Changes in
Interpretation}\label{changes-in-interpretation}

\begin{itemize}
\tightlist
\item
  σ.k unified from ``wave number / acceleration'' to ``vibrational
  frequency''
\item
  δt clarified from ``proper time step'' to ``processing step period''
\item
  ``Time'' is a derived quantity appearing as V{[}t-axis{]}.θ₀
\end{itemize}

\subsubsection{§9.2 Structural Extension}\label{structural-extension}

\begin{itemize}
\tightlist
\item
  δV and δδt added to Particle (v1.0's Particle is included as a subset)
\item
  Mesh structure newly introduced (not present in v1.0)
\end{itemize}

\subsubsection{§9.3 Compatibility with {[}W8{]} Table
A}\label{compatibility-with-w8-table-a}

The 62-particle classification of {[}W8{]} Table A remains valid in
v2.0. However, the Higgs, instead of being counted as a Particle, is
\textbf{counted as an excitation of the Mesh's H\_xyz\_inphase mode} ---
an interpretive update. The total count (61 SM states + 1 graviton = 62)
is unchanged (only the counting interpretation is updated).

\begin{center}\rule{0.5\linewidth}{0.5pt}\end{center}

\subsection{§10 Conclusion}\label{conclusion}

Essence of v2.0:

\begin{itemize}
\tightlist
\item
  \textbf{xyzt are frequencies} (natural interpretation in phase space)
\item
  \textbf{6 axes are completely symmetric} (4 position-type + R + Q)
\item
  \textbf{Transverse (Particle) and longitudinal (Mesh) are separated}
\item
  \textbf{Acceleration is δδt} (change in processing period)
\item
  \textbf{Time is V{[}t{]}.θ₀} (a position phase, advancing per
  particle)
\item
  \textbf{All operations close under int256 integer arithmetic}
\end{itemize}

With this structure, the dynamics rules (D2, separate paper) can
describe:

\begin{itemize}
\tightlist
\item
  Gravitational time dilation, polarization rotation (free motion +
  background coupling)
\item
  Higgs mechanism (coupling to Mesh's H\_xyz\_inphase)
\item
  Gravitational wave emission (excitation of Mesh's H\_xy\_anti)
\item
  Coulomb 2-body problem (inter-particle interaction)
\end{itemize}

all within the same framework.

This paper is a design specification; these dynamical applications are
deferred to separate papers.

\begin{center}\rule{0.5\linewidth}{0.5pt}\end{center}

\subsection{References}\label{references}

\begin{itemize}
\tightlist
\item
  {[}W3{]} Kihara, N. ``Structural Requirements for a Theory of
  Everything.'' DOI: 10.5281/zenodo.19601592 (2026).
\item
  {[}W8{]} Kihara, N. ``Combinatorial Structure of the 6-Dimensional
  Hyperrectangle.'' DOI: 10.5281/zenodo.19762134 (2026).
\item
  {[}W9{]} Kihara, N. ``sine-Gordon Equation --- Foundation of
  Topological Solitons.'' DOI: 10.5281/zenodo.19650966 (2026).
\item
  {[}W10{]} Kihara, N. ``Four Modes of Shape-Invariant Waves.'' DOI:
  10.5281/zenodo.19709798 (2026).
\item
  {[}W11{]} Kihara, N. ``Interactions of Shape-Invariant Waves.'' DOI:
  10.5281/zenodo.19763463 (2026).
\item
  {[}M1{]} Kihara, N. ``Three-Layer Model of Central Projection:
  R--R₁--R₀.'' DOI: 10.5281/zenodo.19691713 (2026).
\item
  {[}TE\_Raxis{]} Kihara, N. ``Reexamination of the 6-Dimensional Sign
  Vector xyztRQ.'' DOI: 10.5281/zenodo.19904714 (2026).
\end{itemize}

\begin{center}\rule{0.5\linewidth}{0.5pt}\end{center}

\textbf{Author}: Noriaki Kihara \textbf{ORCID}:
\href{https://orcid.org/0009-0004-6753-4020}{0009-0004-6753-4020}
\textbf{License}: CC BY 4.0

\end{document}
